% Section 6.7, from lectures/L06/appendix/b_exercises.md.

\section{Uppgifter}
\label{c6:sec:uppgifter}

\begin{aside}
\textbf{Så arbetar du med uppgifterna.} Del~1 och del~2 räknas under lektionen: först på egen
hand, några minuter per uppgift, sedan gemensamt. Del~3 är extrauppgifter att träna vidare på
efter lektionen. De flesta går att räkna i huvudet eller med papper och penna.

\facitnotis{L06}
\end{aside}

% ------------------------------------------------------------------------------------------
\uppgiftsdel{Del 1}{Repetitionsuppgifter}

\uppgift{1.1}{Potensregler}
Förenkla, med positiva exponenter i svaret.
\begin{delar}(3)
  \task $\dfrac{10^6 \cdot 10^{-3}}{10^2}$
  \task $(2 \cdot 10^3)^2$
  \task $\sqrt{4 \times 10^6}$
\end{delar}

\uppgift{1.2}{Standardform i kretsberäkning}
En krets har $U = 3{,}3\,\text{V}$ och $R = 4{,}7\,\text{k}\Omega = 4{,}7 \times 10^3\,\Omega$.
\begin{parts}
  \item Beräkna strömmen $I = \dfrac{U}{R}$ och ange svaret i mA på standardform.
  \item Beräkna effekten $P = \dfrac{U^2}{R}$ och ange svaret i mW.
\end{parts}

\uppgift{1.3}{Kvadratrot}
\begin{parts}
  \item Beräkna $\sqrt{144}$.
  \item Beräkna $\sqrt{0{,}0049}$.
  \item En resistor dissiperar $P = 2\,\text{W}$ vid spänningen $U$, och $R = 200\,\Omega$.
    Beräkna $U$ med hjälp av $P = \dfrac{U^2}{R}$.
\end{parts}

% ------------------------------------------------------------------------------------------
\uppgiftsdel{Del 2}{Nytt stoff}

\uppgift{2.1}{Ren andragradsekvation}
Strömmen $I$ (i ampere) uppfyller
\[
  50I^2 = 8.
\]
\begin{parts}
  \item Lös ekvationen för $I$.
  \item Motivera varför den negativa roten bör förkastas i detta eltekniska sammanhang.
\end{parts}

\uppgift{2.2}{Andragradsekvation med PQ-formeln}
Lös följande ekvationer med PQ-formeln.
\begin{delar}(3)
  \task $x^2 - 5x + 6 = 0$
  \task $x^2 + 3x - 10 = 0$
  \task $x^2 - 6x + 9 = 0$
\end{delar}

\uppgift{2.3}{Andragradsekvation med ABC-formeln}
Lös med ABC-formeln.
\begin{delar}(3)
  \task $2x^2 - 7x + 3 = 0$
  \task $3x^2 + 5x - 2 = 0$
\end{delar}

\uppgift{2.4}{Seriekopplade motstånd}
Två seriekopplade motstånd $R_1$ och $R_2$ uppfyller
\[
  R_1 + R_2 = 13\,\Omega \qquad \text{och} \qquad R_1 \cdot R_2 = 36\,\Omega^2.
\]
\begin{parts}
  \item Använd Vietas formler för att sätta upp en andragradsekvation med $R_1$ som obekant.
  \item Lös ekvationen och ange båda motståndsvärdena.
\end{parts}

% ------------------------------------------------------------------------------------------
\uppgiftsdel{Del 3}{Extrauppgifter}
Uppgifterna nedan är extra träning och görs med fördel på egen hand efter lektionen.

\uppgift{3.1}{Rena andragradsekvationer}
Lös ekvationerna.
\begin{delar}(3)
  \task $x^2 = 49$
  \task $3x^2 = 27$
  \task $x^2 - 5 = 20$
  \task $2x^2 + 8 = 0$
  \task $\dfrac{x^2}{4} = 9$
\end{delar}

\uppgift{3.2}{Faktoriseringsmetoden}
Lös genom att faktorisera och använda nollproduktmetoden.
\begin{delar}(3)
  \task $x^2 - 7x = 0$
  \task $x^2 - 16 = 0$
  \task $x^2 + 5x + 6 = 0$
  \task $2x^2 - 8x = 0$
  \task $x^2 - 4x + 4 = 0$
\end{delar}

\uppgift{3.3}{PQ-formeln}
Lös med PQ-formeln.
\begin{delar}(2)
  \task $x^2 + 2x - 8 = 0$
  \task $x^2 - 8x + 12 = 0$
  \task $x^2 + x - 12 = 0$
  \task $x^2 - 2x - 15 = 0$
\end{delar}

\uppgift{3.4}{Normera först}
Ekvationerna nedan är inte normerade. Dividera först så att $a = 1$ och lös sedan med PQ-formeln.
\begin{delar}(3)
  \task $2x^2 - 10x + 12 = 0$
  \task $3x^2 + 12x - 15 = 0$
  \task $-x^2 + 4x - 3 = 0$
\end{delar}

\uppgift{3.5}{Diskriminanten}
Beräkna diskriminanten $D = b^2 - 4ac$ och ange antalet reella rötter \textbf{utan} att lösa
ekvationen.
\begin{delar}(3)
  \task $x^2 - 6x + 5 = 0$
  \task $x^2 + 4x + 4 = 0$
  \task $x^2 + x + 3 = 0$
  \task $2x^2 - 4x + 2 = 0$
  \task $3x^2 + 2x - 1 = 0$
\end{delar}

\uppgift{3.6}{Kvadratkomplettering}
Lös genom kvadratkomplettering.
\begin{delar}(3)
  \task $x^2 + 4x - 5 = 0$
  \task $x^2 - 10x + 21 = 0$
  \task $x^2 + 6x + 4 = 0$
\end{delar}

\uppgift{3.7}{Vietas formler}
\begin{parts}
  \item En normerad andragradsekvation har rötterna $x_1 = 4$ och $x_2 = -2$. Bestäm $p$ och $q$
    och skriv upp ekvationen.
  \item Kontrollera svaret i a) med PQ-formeln.
  \item Använd Vietas formler för att snabbkontrollera att $x_1 = 4$ och $x_2 = 5$ är rötter till
    $x^2 - 9x + 20 = 0$.
  \item Kan en normerad andragradsekvation ha den dubbla roten $x = 3$? Skriv i så fall upp den.
\end{parts}

\uppgift{3.8}{Ström ur effekt}
Effekten i ett motstånd ges av $P = RI^2$.
\begin{parts}
  \item $P = 12{,}5\,\text{W}$ och $R = 50\,\Omega$. Beräkna $I$.
  \item $P = 0{,}18\,\text{W}$ och $R = 200\,\Omega$. Beräkna $I$ i mA.
  \item En resistor är märkt $470\,\Omega$ och $0{,}5\,\text{W}$. Beräkna den största tillåtna
    strömmen.
  \item Varför förkastas den negativa roten i samtliga deluppgifter?
\end{parts}

\uppgift{3.9}{Spänning ur effekt}
Effekten kan också skrivas $P = \dfrac{U^2}{R}$.
\begin{parts}
  \item $P = 5\,\text{W}$ och $R = 20\,\Omega$. Beräkna $U$.
  \item $P = 0{,}125\,\text{W}$ och $R = 2\,000\,\Omega$. Beräkna $U$.
  \item Ett värmeelement ska avge $P = 1\,000\,\text{W}$ vid $U = 230\,\text{V}$. Vilken
    resistans krävs?
\end{parts}

\uppgift{3.10}{Två motstånd ur summa och produkt}
\begin{parts}
  \item Två motstånd uppfyller $R_1 + R_2 = 20\,\Omega$ och $R_1 R_2 = 96\,\Omega^2$. Bestäm
    motstånden.
  \item Två motstånd uppfyller $R_1 + R_2 = 15\,\Omega$ och $R_1 R_2 = 56\,\Omega^2$. Bestäm
    motstånden.
  \item Kan två motstånd ha summan $10\,\Omega$ och produkten $30\,\Omega^2$? Motivera med
    diskriminanten.
\end{parts}

\uppgift{3.11}{Serie- och parallellresistans samtidigt}
Två motstånd ger tillsammans $10\,\Omega$ seriekopplade och $2{,}4\,\Omega$ parallellkopplade.
\begin{parts}
  \item Visa att $R_1 R_2 = 24\,\Omega^2$.
  \item Sätt upp och lös andragradsekvationen för motstånden.
  \item Kontrollera både serie- och parallellresistansen.
\end{parts}

\uppgift{3.12}{Effekt i en belastning}
Ett okänt motstånd $R$ är seriekopplat med $R_1 = 4\,\Omega$ och matas med $E = 12\,\text{V}$.
Effekten som utvecklas i $R$ ska bli $P = 8\,\text{W}$. Strömmen i kretsen är
$I = \dfrac{E}{R_1 + R}$ och effekten i $R$ är $P = RI^2$.
\begin{parts}
  \item Visa att kravet leder till ekvationen $R^2 - 10R + 16 = 0$.
  \item Lös ekvationen.
  \item Kontrollera \textbf{båda} lösningarna genom att beräkna strömmen och effekten.
  \item Beräkna effekten i $R$ då $R = 4\,\Omega$ och jämför med de två lösningarna.
\end{parts}

% Kept whole on one page: its last part was left alone at the top of the next.
\needspace{16\baselineskip}
\uppgift{3.13}{Sant eller falskt}
Avgör om påståendet är sant eller falskt och motivera kortfattat.
\begin{parts}
  \item Varje andragradsekvation har två reella rötter.
  \item Ekvationen $x^2 = 16$ har lösningen $x = 4$.
  \item PQ-formeln kan användas direkt på $2x^2 + 4x - 6 = 0$.
  \item Om $D = 0$ sammanfaller de båda rötterna.
  \item Om en produkt av två faktorer är noll måste minst en av faktorerna vara noll.
  \item Ekvationen $x^2 + 4 = 0$ saknar reella lösningar.
\end{parts}
